%% ma: L12-B04-0004
%% lop: 12
%% chuong: 1
%% bai: 4
%% dang: 12.4.4
%% loai: TN4
%% muc: TH
%% dap_an: "C"
%% nguon: "Ảnh giáo viên gửi 11/10/2026 (Đề kiểm tra giữa HK1 Toán 12, THPT Việt Đức 2026-2027, mã đề 0103), Phần I Câu 5"
%% kien_thuc: "Bài 4 (đồ thị hàm số bậc ba: cực trị, giao điểm với trục Oy, nhận dạng đồ thị)"
%% trang_thai: chinh-thuc
%% ghi_chu: "Đã duyệt 12/10/2026. Hình TikZ vẽ y = x^3 - 3x + 3 theo các số liệu đọc từ hình gốc (cực đại (-1;5), cực tiểu (1;1), cắt Oy tại (0;3))."
\begin{ex}
Cho hàm số $y=f(x)$ liên tục trên $\mathbb R$ và có đồ thị như hình vẽ.
\begin{center}
\begin{tikzpicture}[x=.9cm,y=.55cm]
  \draw[truc] (-3.4,0)--(5.4,0) node[below]{$x$};
  \draw[truc] (0,-.8)--(0,7.4) node[right]{$y$};
  \draw[thick,domain=-2.2:2.2,samples=80] plot (\x,{\x*\x*\x-3*\x+3});
  \draw[phu] (-1,0)--(-1,5)--(0,5) (1,0)--(1,1)--(0,1);
  \fill (0,3) circle (1.3pt) (-1,5) circle (1.3pt) (1,1) circle (1.3pt);
  \foreach \t in {-3,-2,-1,1,2,3,4,5} \draw (\t,.08)--(\t,-.08) node[below,font=\footnotesize]{$\t$};
  \foreach \t in {1,2,3,5,6} \draw (.05,\t)--(-.05,\t) node[left,font=\footnotesize]{$\t$};
  \node[below left,font=\footnotesize] at (0,0) {$O$};
  \node[right,font=\footnotesize] at (1.9,6.1) {$y=f(x)$};
\end{tikzpicture}
\end{center}
Hỏi $f(x)$ là hàm số nào trong các hàm số dưới đây?
\choice{$y=\dfrac{3x^2-2x+1}{x-1}$}{$y=\dfrac{x-1}{x+2}$}{\True $y=x^3-3x+3$}{$y=x^3-3x^2+3$}
\loigiai{Đồ thị có hai điểm cực trị $(-1;5)$ (cực đại), $(1;1)$ (cực tiểu) và cắt trục tung tại $(0;3)$, là đồ thị hàm bậc ba có $a>0$; do đó loại hai hàm phân thức ở A và B. Với $y=x^3-3x^2+3$: $y(-1)=-1\ne5$ nên loại D. Với $y=x^3-3x+3$: $y(-1)=5$, $y(1)=1$, $y(0)=3$ khớp đồ thị. Vậy chọn C.}
\end{ex}
